\section{Historical Controversies Without Devotional Panic}

Guadalupe has generated opposing apologetics because it bears extraordinary religious and national weight. One side treats every silence, brushstroke, or textual parallel as disproof. Another treats every annal, papal visit, eye enlargement, or conversion story as cumulative scientific demonstration. Both approaches ask evidence to perform outside its competence.

\subsection{The date of the detailed story}

The principal difficulty is not that nothing Guadalupan existed before 1648. The image and cult are earlier. The difficulty is that the complete Juan Diego narrative enters the located printed record in 1648--1649, more than a century after the reported event, while central earlier writers are silent.

Several historical models remain possible:

\begin{enumerate}
  \item a substantially historical 1531 event transmitted orally and through lost or unlocated writings until the seventeenth-century publications;
  \item an early cult and image around which a fuller narrative developed through memory, preaching, theater, theology, and community identity;
  \item a sixteenth-century narrative substantially shaped or rewritten in the seventeenth century;
  \item a principally seventeenth-century hagiographic construction projected back to the cult's origin; or
  \item a mixed history in which real persons, early devotion, later literary formation, and theological interpretation cannot be divided by one clean line.
\end{enumerate}

The extant evidence does not compel every scholar to the same model. Catholic authority has received the tradition and canonized Juan Diego; historians still must state the date and genre of the witnesses they use. A papal homily can authoritatively interpret the received tradition without becoming a newly discovered sixteenth-century archive.

\subsection{The 1556 controversy and the alleged painter Marcos}

Bustamante's sermon complained that the image had been painted by an Indigenous artist called Marcos and that miracle preaching risked misleading the faithful. Archbishop Montúfar's inquiry gathered testimony in a conflict that also involved Franciscan and episcopal approaches to Indigenous evangelization. The source is early and indispensable.

It does not, standing alone, identify a surviving contract, date of painting, workshop, complete technique, or the exact present layer associated with a named artist. An opponent's assertion may preserve true local knowledge; it may also function rhetorically. Responsible art history places the allegation beside examination of the object, copies, engravings, pigments, iconographic sources, and conservation history. Responsible devotion does not suppress it.

The theological consequence is modest. If human hands contributed to the image, Catholic veneration is not defeated. If some feature remains materially unexplained, supernatural origin does not follow automatically. The Church's actual reception should not be replaced by an all-or-nothing authorship contest.

\subsection{The 1666 \emph{Informaciones}}

Ecclesiastical testimony gathered in 1665--1666 is often called a canonical investigation proving the apparition. It is indeed an important formal inquiry into the tradition and cult, undertaken as the cathedral chapter sought papal approval of a feast and proper office. Timothy Matovina's critical review counts twenty witnesses: twelve Spanish or criollo witnesses questioned in Mexico City and eight residents of Cuauhtitlan, seven Indigenous persons speaking through an interpreter and one mestizo. Witnesses transmitted local and family memories, and the proceedings entered later arguments for liturgical recognition.

The dates and interrogation impose limits. A person old enough to testify in 1666 was ordinarily not an adult eyewitness in 1531. The questions themselves summarized the apparition account, and many responses tracked them closely; the 1648 and 1649 books already existed. Details volunteered beyond the interrogatory can have a different evidentiary value, especially family chains and local memories, but they still require source criticism. The proceedings are not equivalent to depositions recorded during the reported week, nor were they a modern DDF decision with an operative supernatural-origin formula. Calling them “canonical” describes form, not contemporaneity or a conclusion beyond the actual text.

\subsection{Juan Diego's historicity and canonization}

The late narrative record produced public dispute during the twentieth-century cause. Supporters assembled annals, oral transmission, cult, linguistic arguments, and documentary studies. Critics emphasized Zumárraga's silence, late texts, uncertain manuscripts, and the development of the story. John Paul II beatified Juan Diego in 1990 and canonized him in 2002.

Canonization is a decisive ecclesial act regarding the saint: the Church authoritatively proposes Juan Diego's heavenly glory, holiness, and public cult. It is not an inspired critical edition of the \emph{Nican mopohua}. The homily itself centers his humble discipleship, inculturated reception of Christ, family and social virtue, and service as Mary's messenger. Catholics owe respect to the saint without being required to claim that every supporting document is contemporary or every historical objection frivolous.

Historians who are Catholic may differ over how the person is recovered from the sources. They may not misstate the canonization as if it did not occur. Devotees may confidently invoke Saint Juan Diego. They may not convert canonization into a verdict on pigment, Nahuatl authorship, or an alleged astronomical map.

\subsection{The image's material history}

The image has been displayed, handled, mounted, framed, copied, crowned, photographed, protected, and exposed across centuries. A juridical file opened in 1820 and continued through 1823 gathered later testimony about an accidental \emph{agua fuerte} spill, probably around 1784 or 1785 but not securely dated. On 14 November 1921 a bomb exploded below the altar. A contemporary newspaper account described damage to the altar and candlestick, a bent bronze crucifix, and an intact image, frame, and protective glass. The object's survival entered devotional memory. Historical reproductions and technical discussions also indicate changes or interventions in features such as the crown, rays, moon, angel, gold, and surface.

The project did not perform imaging or sampling. Published technical discussion is fragmented, access to the original is controlled, and apologetic summaries regularly combine unlike observers into a fictional unanimous “scientific study.” A conservation conclusion should identify examiner, date, access, instruments, wavelengths or samples, image registration, chain of custody, raw data, condition history, uncertainty, and peer review. Without these, confidence must be proportionate.

\begin{fourstable}{Claim type}{What is secure}{What is disputed or absent}{Proportionate conclusion}
Early cult & Chapel, image, pilgrims, offerings, and controversy by mid-1550s & Exact origin story in those early witnesses & Early cult is historical; full narrative is not thereby contemporaneously documented.\\
Full narrative & Complete Spanish 1648 and Nahuatl 1649 witnesses & Earlier archetype, authorship, degree of dependence and oral continuity & Read the received narrative fully while dating the extant texts honestly.\\
Juan Diego & Canonized saint; strong ecclesial reception; later narrative and supporting traditions & Exact reconstruction of every biographical detail & Invoke the saint; do not turn the cause into universal historical inerrancy.\\
Sacred image & Venerated material object with stable core iconography and long history & Maker, technique, layers, restorations, complete material mechanism & Conserve and study; neither declare fraud nor miracle beyond evidence.\\
Ecclesial status & Feast, shrine, papal teaching, patronage, canonization & Located modern event-origin decree & Call the reception exceptional and the juridic object exact.\\
\end{fourstable}

\subsection{Fruit, identity, and truth}

The devotion has produced prayer, conversion, art, solidarity, vocations, care for migrants, defense of vulnerable life, and a shared language across the Americas. It has also been commercialized, nationalized, and enlisted by incompatible political movements. Good fruit strongly matters in ecclesial discernment, but it does not retroactively date a manuscript. Abuse does not invalidate every authentic fruit.

Truthfulness is itself a fruit. A devotion that cannot survive a footnote about 1649 has been built on the wrong foundation. A historical method that cannot recognize centuries of holiness and communal meaning because it lacks a modern event dossier has also narrowed reality. The mature reader does not fear a mixed evidentiary landscape.

\begin{studybox}{Rule for disputed claims}
Name the exact proposition, identify the earliest and strongest witness, distinguish event from cult and interpretation, report serious contrary evidence, and conclude no more than the evidence permits. Guadalupe loses nothing when an unsupported marvel is surrendered; its deepest evangelical meaning becomes clearer.
\end{studybox}
